%===============================================================================
%  APPENDIX: NOMENCLATURE
%===============================================================================
\section{Nomenclature and Mathematical Notation}
\label{app:nomenclature}
%===============================================================================

In this section, we provide a comprehensive list of the mathematical notations used throughout our work, organized by category.
This serves as a centralized reference to ensure clarity and consistency.

%-------------------------------------------------------------------------------
\subsection{General Variables and Data Structures}
\label{app:nomenclature:general}
%-------------------------------------------------------------------------------

We begin by defining the fundamental variables and data structures used to represent observed and latent spaces.

\begin{table}[H]
\begin{tabularx}{\textwidth}{l X}
\toprule
\textbf{Symbol} & \textbf{Description} \\
\midrule
$\realnumbers$ & The set of real numbers. \\
$\variables$ & A vector of observed variables (e.g., gene expression). \\
$\latentvars$ & A vector of general latent variables in the learned representation space. \\
$\noiseterms$ & A vector of general, unspecified noise terms. \\
$\datamatrix$ & A matrix representing a set of observed data samples. \\
$\identitymatrix$ & The identity matrix. \\
\bottomrule
\end{tabularx}
\end{table}

%-------------------------------------------------------------------------------
\subsection{Causal Models and Graphs}
\label{app:nomenclature:causal}
%-------------------------------------------------------------------------------

Next, we define the components of the Structural Causal Model (SCM) that describes the underlying data-generating process.

\begin{table}[H]
\begin{tabularx}{\textwidth}{l X}
\toprule
\textbf{Symbol} & \textbf{Description} \\
\midrule
$\graph$ & A \gls{dag} representing the causal structure. \\
$\graphmatrix$ & The weighted adjacency matrix of the causal graph $\graph$. \\
$\adjacencymatrix$ & A general adjacency matrix. \\
$\weightmatrix$ & A general weight matrix. \\
$\parentnodes(\cdot)$ & A function that returns the parent nodes of a given node in $\graph$. \\
$\transitiveclosure{\graph}$ & The transitive closure of the graph $\graph$. \\
$\latentcausalvars$ & The true, unobserved latent causal variables in the \gls{scm}. \\
$\latentcausalvar$ & A single latent causal variable. \\
$\exonoise$ & A vector of exogenous noise terms for the \gls{scm}. \\
$\structuralfunc_i$ & The structural causal function for a single variable $\latentcausalvar_i$. \\
$\tr(\cdot)$ & The trace function for a matrix. \\
\bottomrule
\end{tabularx}
\end{table}

%------------------------------------------------------------------------------
\subsection{Model Functions and Representations}
\label{app:nomenclature:model_functions}
%------------------------------------------------------------------------------

We denote the core functions of our model, which learn to map between the observed and latent spaces, as follows.

\begin{table}[H]
\begin{tabularx}{\textwidth}{l X}
\toprule
\textbf{Symbol} & \textbf{Description} \\
\midrule
$\truemixingfunc$ & The true, unknown data-generating (mixing) function. \\
$\decoder$ & The learned decoder network. \\
$\idealencoder$ & The ideal, oracle encoder function. \\
$\practicalencoder$ & The practical, parameterized encoder network. \\
$\learnedrepresentation$ & The learned latent representation, output of the encoder ($\practicalencoder(\variables)$). \\
$\predictedvariable$ & A single data sample generated by the model's decoder. \\
$\predicteddatamatrix$ & A matrix representing a batch of data samples generated by the model's decoder. \\
$\composedfunc$ & The composed function $\practicalencoder \circ \truemixingfunc$. \\
$\latentideal$ & The true latent representation produced by the $\idealencoder$. \\
$\latentchangeideal$ & The ideal, single-node change in the latent space. \\
$\latentchangeapprox$ & The approximated (dense) change in the latent space. \\
$\normaldist$ & The normal (Gaussian) distribution. \\
%$\auxvar$ & An auxiliary conditioning variable (Khemakhem et al., 2020). \\
%$a$ & The composed function, $\practicalencoder \circ \truemixingfunc$. \\
%$p$ & The degree of a polynomial mixing function. \\
\bottomrule
\end{tabularx}
\end{table}

%-------------------------------------------------------------------------------
\subsection{Intervention-Specific Variables}
\label{app:nomenclature:interventions}
%-------------------------------------------------------------------------------

Here, we list the variables specifically related to the modeling of interventions.

\begin{table}[H]
\begin{tabularx}{\textwidth}{l X}
\toprule
\textbf{Symbol} & \textbf{Description} \\
\midrule
$\interventionencoder$ & The intervention encoder network. \\
$\vi$ & The one-hot vector specifying the target of an intervention. \\
$\conditionvector$ & A general condition vector input to an intervention network. \\
$\maskvector$ & The mask vector specifying an intervention's target(s). \\
$\interventionstrength$ & The scalar strength of a latent intervention. \\
$\zobs$ & Latents from an observational (control) sample. \\
$\zint$ & Latents from an interventional sample. \\
$\interventionaldist$ & The interventional distribution over latent variables. \\
$\interventionset$ & The set of intervened variables in environment $k$. \\
$\interventiontargets$ & The targets of an intervention for a variable $i$. \\
$\auxvar$ & An auxiliary variable providing conditioning information (e.g., environment index). \\
\bottomrule
\end{tabularx}
\end{table}

%-------------------------------------------------------------------------------
\subsection{Identifiability and Transformations}
\label{app:nomenclature:identifiability}
%-------------------------------------------------------------------------------

We define symbols used in the context of identifiability proofs and the resulting transformation classes.

\begin{table}[H]
\begin{tabularx}{\textwidth}{l X}
\toprule
\textbf{Symbol} & \textbf{Description} \\
\midrule
$\simT$ & Equivalence class for a transformation T. \\
$\simC$ & Equivalence class for a transformation C. \\
$\affinetransform$ & The matrix component of an affine transformation. \\
$\affineoffset$ & The vector component of an affine transformation. \\
$\scalefactor$ & A scalar factor in a scaling transformation. \\
$\shiftoffset$ & An offset in a shift transformation. \\
\bottomrule
\end{tabularx}
\end{table}

%-------------------------------------------------------------------------------
\subsection{Probability, Losses, and Functions}
\label{app:nomenclature:probability_losses}
%-------------------------------------------------------------------------------

We define the probability functions, loss functions, and mathematical operators used to train our models and enforce constraints.

\begin{table}[H]
\begin{tabularx}{\textwidth}{l X}
\toprule
\textbf{Symbol} & \textbf{Description} \\
\midrule
$\prob{\cdot}$ & \gls{pmf} (for discrete variables). \\
$\pdf{\cdot}$ & \gls{pdf} (for continuous variables). \\
$\cdf{\cdot}$ & \gls{cdf}. \\
$\pdata$ & The true, underlying data-generating distribution (\gls{pdf}). \\
$\pmodel$ & The learned distribution of the model (\gls{pdf}). \\
$\E[\cdot]$ & The expectation function. \\
$\dagmaloss$ & Loss function related to the causal graph structure. \\
$\mmdloss$ & Loss function based on the \gls{mmd}. \\
$\acyclicityfunc$ & A function used to enforce graph acyclicity. \\
$\scorefunc$ & A score function for graph discovery. \\
$\mmd(\cdot, \cdot)$ & The \gls{mmd} function, which computes the distance between two distributions. \\
$\mkmmd(\cdot, \cdot)$ & The \gls{mmd} function, which computes the distance between two distributions. \\
$\kernel(\cdot, \cdot)$ & A kernel function used in \gls{mmd} computation. \\
\bottomrule
\end{tabularx}
\end{table}

%-------------------------------------------------------------------------------
\subsection{Hyperparameters}
\label{app:nomenclature:hyperparams}
%-------------------------------------------------------------------------------

We list the key hyperparameters that control our model's training and behavior.

\begin{table}[H]
\begin{tabularx}{\textwidth}{l X}
\toprule
\textbf{Symbol} & \textbf{Description} \\
\midrule
$\scorelambda$ & Regularization hyperparameter for a graph score. \\
$\mmdsigmaMMD$ & The bandwidth hyperparameter for the Gaussian RBF kernel. \\
$\mkmmdweights$ & The weights for the convex combination in multi-kernel MMD. \\
$\numkernels$ & The number of kernels used in the \gls{mmd} calculation. \\
$\numgraphs$ & The number of graphs. \\
$\numsamplesobs$ & The number of samples in a batch of observed data. \\
$\numsamplespred$ & The number of samples in a batch of predicted data. \\
\bottomrule
\end{tabularx}
\end{table}

%-------------------------------------------------------------------------------
\subsection{Downstream Analysis Metrics}
\label{app:nomenclature:downstream_metrics}
%-------------------------------------------------------------------------------

We define variables used in the downstream evaluation tasks, such as the genetic interaction analysis.

\begin{table}[H]
\begin{tabularx}{\textwidth}{l X}
\toprule
\textbf{Symbol} & \textbf{Description} \\
\midrule
$\effectvector$ & The perturbation effect vector (mean of perturbed minus mean of control). \\
$\effectvectornaive$ & The naive additive effect vector, defined as $\effectvector A + \effectvector B$. \\
$\ltwonorm{\cdot}$ & The L2 norm of a vector. \\
$\cossim{\cdot}{\cdot}$ & The cosine similarity between two vectors. \\
$\meanexpression$ & The mean expression profile of a cell population. \\
$\numcells$ & The number of cells for a given perturbation condition. \\
$\predictionfunc_M$ & The prediction function for a given model $M$. \\
$\synergyscore$ & The score used to measure synergy and suppression. \\
$\neomorphismscore$ & The score used to measure neomorphism. \\
$\redundancyscore$ & The score used to measure redundancy. \\
$\epistasiscore$ & The score used to measure epistasis. \\
\bottomrule
\end{tabularx}
\end{table}

%-------------------------------------------------------------------------------
\subsection{Dimensionality and Spaces}
\label{app:nomenclature:dims}
%-------------------------------------------------------------------------------

Finally, we list symbols for dimensionality and the properties of the data spaces.

\begin{table}[H]
\begin{tabularx}{\textwidth}{l X}
\toprule
\textbf{Symbol} & \textbf{Description} \\
\midrule
$\observeddim$ & The dimensionality of the observed data space. \\
$\latentdim$ & The dimensionality of the latent space. \\
$\numpathways$ & The number of learned pathways, equivalent to $\latentdim$. \\
$\polydegree$ & The degree of a polynomial mixing function. \\
$\latentsupport$ & The support of the latent distribution. \\
$\observsupport$ & The support of the observed distribution. \\
$\numgenes$ & The number of genes in the expression profile, equivalent to $\observeddim$. \\
\bottomrule
\end{tabularx}
\end{table}

%-------------------------------------------------------------------------------
\subsection{Acronyms and Abbreviations}
\label{app:nomenclature:acronyms}
%-------------------------------------------------------------------------------

\printglossary[type=\acronymtype,title={List of Acronyms}]
